\documentclass[11pt,letterpaper]{article}
% This source is intentionally restricted to pdfLaTeX.
\usepackage[T1]{fontenc}
\usepackage[utf8]{inputenc}
\usepackage{textcomp}
\usepackage{lmodern}
\usepackage[margin=0.82in,top=0.94in,bottom=0.91in,headheight=18pt]{geometry}
\usepackage{amsmath,amssymb}
\usepackage[table]{xcolor}
\definecolor{DeepNavy}{HTML}{14283F}
\definecolor{Teal}{HTML}{0A747C}
\definecolor{Gold}{HTML}{C78F36}
\definecolor{Slate}{HTML}{536477}
\definecolor{Pale}{HTML}{F4F7FA}
\definecolor{CodeEdge}{HTML}{D4DFE9}
\usepackage{microtype}
\usepackage{parskip}
\setlength{\parskip}{0.47em}
\usepackage{titlesec}
\titleformat{\section}{\normalfont\Large\bfseries\color{DeepNavy}}{\thesection}{0.62em}{}
\titlespacing*{\section}{0pt}{2.0ex plus .7ex minus .2ex}{0.75ex}
\titleformat{\subsection}{\normalfont\large\bfseries\color{Teal}}{\thesubsection}{0.6em}{}
\usepackage{longtable,booktabs,array,calc}
\renewcommand{\arraystretch}{1.22}
\AtBeginEnvironment{longtable}{\small}
\usepackage{etoolbox}
\makeatletter
\patchcmd\longtable{\par}{\if@noskipsec\mbox{}\fi\par}{}{}
\makeatother
\usepackage{listings}
\IfFileExists{upquote.sty}{\usepackage{upquote}}{}
\lstdefinestyle{suitecode}{
  basicstyle=\ttfamily\footnotesize,
  breaklines=true,breakatwhitespace=false,
  columns=fullflexible,keepspaces=true,showstringspaces=false,
  frame=single,rulecolor=\color{CodeEdge},
  backgroundcolor=\color{Pale},
  keywordstyle=\color{Teal}\bfseries,
  commentstyle=\color{Slate}\itshape,
  stringstyle=\color{DeepNavy},
  xleftmargin=5pt,xrightmargin=5pt,
  framexleftmargin=5pt,framexrightmargin=5pt,
  aboveskip=0.85em,belowskip=0.95em,
  tabsize=2
}
\lstset{style=suitecode}
\usepackage{enumitem}
\setlist[itemize]{leftmargin=1.6em,itemsep=2pt,topsep=3pt}
\setlist[enumerate]{leftmargin=1.8em,itemsep=2pt,topsep=3pt}
\providecommand{\tightlist}{\setlength{\itemsep}{0pt}\setlength{\parskip}{0pt}}
\newcommand{\passthrough}[1]{#1}
\usepackage{xurl}
\usepackage[colorlinks=true,linkcolor=Teal,urlcolor=Teal,citecolor=Teal,pdfborder={0 0 0}]{hyperref}
\urlstyle{same}
\usepackage{fancyhdr,lastpage}
\pagestyle{fancy}
\fancyhf{}
\fancyhead[L]{\footnotesize\sffamily\color{Slate} ECMA-262:2025 \textbullet{} Engineering Documentation}
\fancyhead[R]{\footnotesize\sffamily\color{Slate} 16th Edition}
\fancyfoot[L]{\footnotesize\sffamily\color{Slate} Independent technical companion}
\fancyfoot[R]{\footnotesize\sffamily\color{Slate}\thepage\ / \pageref{LastPage}}
\renewcommand{\headrulewidth}{0.35pt}
\renewcommand{\footrulewidth}{0pt}
\setlength{\emergencystretch}{3em}
\setcounter{tocdepth}{2}
\setcounter{secnumdepth}{2}
\begin{document}
\begin{titlepage}
\thispagestyle{empty}
\vspace*{2.8cm}
{\color{Gold}\rule{18mm}{3pt}}\par\vspace{18mm}
{\sffamily\Large\bfseries\color{Teal} SOFTWARE ENGINEERING DOCUMENTATION SUITE}\par
\vspace{8mm}
{\sffamily\huge\bfseries\color{DeepNavy} ECMAScript 2025 }\par
\vspace{4mm}
{\sffamily\fontsize{24}{29}\selectfont\bfseries\color{DeepNavy} Language Technical Reference }\par
\vspace{10mm}
{\large\color{Slate} Types, syntax, semantics and 2025 additions }\par
\vspace{16mm}
{\color{Teal}\rule{\textwidth}{0.8pt}}\par
\vfill
{\sffamily\color{Slate}\large ECMA-262:2025 \textbullet{} 16th Edition \textbullet{} June 2025}\par
\vspace{3mm}
{\sffamily\color{Slate} October 2026 \textbullet{} Document 1 of 8}\par
\vspace{3mm}
{\small\color{Slate} Original engineering companion. The uploaded ECMA-262:2025 specification remains the authoritative source.}\par
\end{titlepage}
\pagenumbering{arabic}
\tableofcontents
\clearpage
\section{Scope, authority, and navigation}\label{scope-authority-and-navigation}

\textbf{Edition baseline.} This reference interprets ECMA-262, \emph{ECMAScript 2025 Language Specification}, 16th edition (June 2025), clauses 1--29 and Annexes A--F. It is an original engineering companion, not a replacement for the standard. References in parentheses point to the supplied edition. Conformance obligations attach to implementations, not to arbitrary application coding-style rules (clause 2).

\textbf{Classification.} \textbf{{[}SPEC{]}} means a summarized language or implementation requirement; \textbf{{[}HOST{]}} denotes behavior supplied by a browser or another host; \textbf{{[}GUIDE{]}} is an engineering recommendation not imposed by ECMA-262. Annex B is normative for web-browser hosts but normative-optional for other hosts. Annexes A and C--F are informative. Legacy and normative-optional subclauses require careful reading of the annotations in clause 2.

\begin{longtable}[]{@{}
  >{\raggedright\arraybackslash}p{(\columnwidth - 4\tabcolsep) * \real{0.3333}}
  >{\raggedright\arraybackslash}p{(\columnwidth - 4\tabcolsep) * \real{0.3333}}
  >{\raggedright\arraybackslash}p{(\columnwidth - 4\tabcolsep) * \real{0.3333}}@{}}
\toprule\noalign{}
\begin{minipage}[b]{\linewidth}\raggedright
Specification clauses
\end{minipage} & \begin{minipage}[b]{\linewidth}\raggedright
Technical subject
\end{minipage} & \begin{minipage}[b]{\linewidth}\raggedright
Common engineering question
\end{minipage} \\
\midrule\noalign{}
\endhead
\bottomrule\noalign{}
\endlastfoot
5--8 & Grammar, types, abstract operations, static semantics & What algorithm determines observable behavior? \\
9--10 & Contexts, realms, ordinary/exotic objects & Where are bindings and prototype invariants enforced? \\
11--16 & Source, syntax, scripts, modules & Is this valid syntax and what does it evaluate to? \\
17--19 & Extensions, built-ins, global object & Which facilities belong to the language? \\
20--28 & Intrinsics and standardized APIs & What does a built-in method promise? \\
29 & Shared-memory execution model & What outcomes are allowed for concurrent agents? \\
Annex B & Legacy web-browser compatibility & May new portable code depend on this behavior? \\
\end{longtable}

\section{Types, values, and equality}\label{types-values-and-equality}

\textbf{{[}SPEC{]}} ECMAScript language types are Undefined, Null, Boolean, String, Symbol, Number, BigInt, and Object (6.1). The specification additionally defines \emph{specification types} such as Reference Records and Completion Records (6.2); these are devices for defining evaluation and are not general-purpose values created by application code. A Number is a binary64 floating-point value with \passthrough{\lstinline!NaN!}, infinities, and signed zero. BigInt holds integral values with different operator constraints.

\begin{longtable}[]{@{}
  >{\raggedright\arraybackslash}p{(\columnwidth - 4\tabcolsep) * \real{0.3333}}
  >{\raggedright\arraybackslash}p{(\columnwidth - 4\tabcolsep) * \real{0.3333}}
  >{\raggedright\arraybackslash}p{(\columnwidth - 4\tabcolsep) * \real{0.3333}}@{}}
\toprule\noalign{}
\begin{minipage}[b]{\linewidth}\raggedright
Operation
\end{minipage} & \begin{minipage}[b]{\linewidth}\raggedright
Rule of thumb
\end{minipage} & \begin{minipage}[b]{\linewidth}\raggedright
Important boundary
\end{minipage} \\
\midrule\noalign{}
\endhead
\bottomrule\noalign{}
\endlastfoot
\passthrough{\lstinline!===!} & Strict equality; does not coerce distinct types & \passthrough{\lstinline"NaN !== NaN"}; \passthrough{\lstinline!0 === -0!} \\
\passthrough{\lstinline!Object.is!} & SameValue comparison & distinguishes positive and negative zero; treats NaNs alike \\
\passthrough{\lstinline!==!} & Abstract equality with conversions & surprising when operands have distinct types \\
\passthrough{\lstinline!Map!} / \passthrough{\lstinline!Set!} key equivalence & SameValueZero & NaNs equal, zeros not distinguished \\
\passthrough{\lstinline!typeof null!} & returns \passthrough{\lstinline!"object"!} & historical behavior; test null explicitly \\
\passthrough{\lstinline!Number(1n)!} & explicit possible precision-loss conversion & avoid for integers outside safe range \\
\end{longtable}

\begin{lstlisting}
console.assert(0 === -0);
console.assert(!Object.is(0, -0));
console.assert(Object.is(NaN, NaN));
console.assert(new Set([NaN, NaN]).size === 1);
console.assert(typeof null === 'object');
console.assert(9007199254740993n > 9007199254740992n);
\end{lstlisting}

\textbf{{[}GUIDE{]}} Use explicit contracts at boundaries: validate numeric ranges, do not mix BigInt and Number arithmetic unintentionally, and select equality to match the domain's invariants.

\section{Conversion and abstract operations}\label{conversion-and-abstract-operations}

\textbf{{[}SPEC{]}} Clause 7 defines reusable abstract operations such as \passthrough{\lstinline!ToPrimitive!}, \passthrough{\lstinline!ToBoolean!}, \passthrough{\lstinline!ToNumber!}, \passthrough{\lstinline!ToNumeric!}, \passthrough{\lstinline!ToString!}, \passthrough{\lstinline!ToPropertyKey!}, \passthrough{\lstinline!Get!}, \passthrough{\lstinline!Set!}, and \passthrough{\lstinline!Call!}. Operators in clauses 13--14 invoke these definitions. \passthrough{\lstinline!ToBoolean!} treats \passthrough{\lstinline!false!}, zero, \passthrough{\lstinline!-0!}, \passthrough{\lstinline!0n!}, \passthrough{\lstinline!NaN!}, empty string, \passthrough{\lstinline!null!}, and \passthrough{\lstinline!undefined!} as falsy; every ordinary object is truthy. Converting an object to a primitive can invoke user-provided methods and cause side effects.

\begin{lstlisting}
const value = { valueOf() { console.log('conversion'); return 4; } };
console.assert(value + 2 === 6); // conversion calls user code
console.assert(Boolean([]) === true);
console.assert('5' + 1 === '51');
console.assert(Number('5') + 1 === 6);
\end{lstlisting}

\textbf{{[}GUIDE{]}} When reading abstract algorithms, track evaluation order, every abrupt completion, the exact abstract operation invoked, and whether it can call user code (\passthrough{\lstinline!valueOf!}, getters, proxies, iteration methods).

\section{Lexical and syntactic structure}\label{lexical-and-syntactic-structure}

\textbf{{[}SPEC{]}} Source text, goal symbols, Unicode handling, hashbang comments, lexical tokens, ASI, and regular-expression lexical decisions are covered in 11--12. Expressions and operators are in clause 13; statements and declarations are in clause 14. Grammar has context parameters (such as \passthrough{\lstinline!Yield!}, \passthrough{\lstinline!Await!}, \passthrough{\lstinline!In!}) and \emph{early errors}, which are syntax restrictions checked before execution. Automatic semicolon insertion is a defined parser rule, not universal permission to omit semicolons.

\begin{lstlisting}
function misleading() {
  return
  { success: true }; // ASI makes the return value undefined
}
console.assert(misleading() === undefined);
\end{lstlisting}

\begin{longtable}[]{@{}
  >{\raggedright\arraybackslash}p{(\columnwidth - 4\tabcolsep) * \real{0.3333}}
  >{\raggedright\arraybackslash}p{(\columnwidth - 4\tabcolsep) * \real{0.3333}}
  >{\raggedright\arraybackslash}p{(\columnwidth - 4\tabcolsep) * \real{0.3333}}@{}}
\toprule\noalign{}
\begin{minipage}[b]{\linewidth}\raggedright
Binding form
\end{minipage} & \begin{minipage}[b]{\linewidth}\raggedright
Scope
\end{minipage} & \begin{minipage}[b]{\linewidth}\raggedright
Initialization / notes
\end{minipage} \\
\midrule\noalign{}
\endhead
\bottomrule\noalign{}
\endlastfoot
\passthrough{\lstinline!var!} & variable environment / function or script semantics & historically hoisted; avoid in new code \\
\passthrough{\lstinline!let!} & lexical scope & uninitialized until declaration is evaluated \\
\passthrough{\lstinline!const!} & lexical scope & initializer required in ordinary declarations; binding immutable \\
\passthrough{\lstinline!class!} declaration & lexical scope & subject to initialization constraints \\
\passthrough{\lstinline!function!} declaration & context-sensitive semantics & inspect Annex B for legacy block-function deviations \\
\end{longtable}

\section{Functions, closures, classes, and objects}\label{functions-closures-classes-and-objects}

\textbf{{[}SPEC{]}} Function creation, argument binding, arrow-function lexical \passthrough{\lstinline!this!}, generators, async functions, private class elements, and class initialization are defined across clauses 10 and 15. Ordinary property access traverses prototypes; property descriptors and internal methods are central to clauses 10 and 20. A proxy may intercept internal methods but must honor specified invariants. Class private names are not ordinary string-named properties.

\begin{lstlisting}
class Ledger {
  #balance = 0;
  deposit(n) {
    if (!Number.isFinite(n) || n <= 0) throw new RangeError('amount');
    this.#balance += n;
    return this.#balance;
  }
  get balance() { return this.#balance; }
}
const ledger = new Ledger();
console.assert(ledger.deposit(2) === 2);
\end{lstlisting}

\textbf{{[}GUIDE{]}} Do not equate \passthrough{\lstinline!Object.freeze!} with recursive immutability. Review getters, setters, property descriptors, prototype changes, object spread, and proxy traps when designing security-sensitive APIs.

\section{Scripts, modules, and import attributes}\label{scripts-modules-and-import-attributes}

\textbf{{[}SPEC{]}} Clause 16 distinguishes scripts from modules and defines linking, evaluation, module requests, and import/export syntax. Modules have strict-mode semantics and live bindings. Module loading, resolution, and certain import attributes involve host-defined behavior: having grammar in ECMA-262 does not imply an arbitrary host will load a named file or URL. ECMAScript 2025 standardizes import attributes and JSON modules (16.2).

\begin{lstlisting}
// Host support is required to resolve and load this module.
import settings from './settings.json' with { type: 'json' };
export const retries = settings.retries;
\end{lstlisting}

\textbf{{[}HOST{]}} Browser and Node.js resolution, file permissions, and module-specifier conventions are outside the complete core language contract. \textbf{{[}GUIDE{]}} Make portability an explicit test target rather than assuming all engines and hosts support the same loading strategy.

\section{Built-in object taxonomy}\label{built-in-object-taxonomy}

\begin{longtable}[]{@{}
  >{\raggedright\arraybackslash}p{(\columnwidth - 4\tabcolsep) * \real{0.3333}}
  >{\raggedright\arraybackslash}p{(\columnwidth - 4\tabcolsep) * \real{0.3333}}
  >{\raggedright\arraybackslash}p{(\columnwidth - 4\tabcolsep) * \real{0.3333}}@{}}
\toprule\noalign{}
\begin{minipage}[b]{\linewidth}\raggedright
Domain
\end{minipage} & \begin{minipage}[b]{\linewidth}\raggedright
Principal clauses
\end{minipage} & \begin{minipage}[b]{\linewidth}\raggedright
Representative APIs
\end{minipage} \\
\midrule\noalign{}
\endhead
\bottomrule\noalign{}
\endlastfoot
Fundamental objects & 20 & Object, Function, Symbol, Error \\
Numbers and time & 21 & Number, BigInt, Math, Date \\
Text processing & 22 & String, RegExp \\
Indexed collections & 23 & Array, TypedArray \\
Keyed collections & 24 & Map, Set, WeakMap, WeakSet \\
Structured data & 25 & ArrayBuffer, SharedArrayBuffer, DataView, Atomics, JSON \\
Memory management & 26 & WeakRef, FinalizationRegistry \\
Control abstraction & 27 & Iterator, Promise, generators, async functions \\
Reflection & 28 & Reflect, Proxy, module namespace objects \\
\end{longtable}

\textbf{{[}GUIDE{]}} Select structures by operation: \passthrough{\lstinline!Map!} for arbitrary keys, \passthrough{\lstinline!Set!} for membership and algebra, typed arrays for binary data, ordinary objects for named records, promises for one-shot asynchronous results. Garbage-collection timing must not be used for correctness decisions.

\section{ECMAScript 2025 change register}\label{ecmascript-2025-change-register}

The edition introduction explicitly identifies the following 2025 additions. The examples below demonstrate intended use; deployment still requires an implementation with relevant feature support.

\begin{longtable}[]{@{}
  >{\raggedright\arraybackslash}p{(\columnwidth - 4\tabcolsep) * \real{0.3333}}
  >{\raggedright\arraybackslash}p{(\columnwidth - 4\tabcolsep) * \real{0.3333}}
  >{\raggedright\arraybackslash}p{(\columnwidth - 4\tabcolsep) * \real{0.3333}}@{}}
\toprule\noalign{}
\begin{minipage}[b]{\linewidth}\raggedright
Addition
\end{minipage} & \begin{minipage}[b]{\linewidth}\raggedright
Primary reference
\end{minipage} & \begin{minipage}[b]{\linewidth}\raggedright
Minimal engineering use
\end{minipage} \\
\midrule\noalign{}
\endhead
\bottomrule\noalign{}
\endlastfoot
Iterator global and helpers & 27.1 & lazy transformation and consumption \\
Set operations & 24.2.4 & intersection, union, difference, relational predicates \\
JSON modules/import attributes & 16.2 & type-annotated JSON imports \\
\passthrough{\lstinline!RegExp.escape!} & 22.2.5.1 & safely embed a literal string in pattern text \\
Inline regexp modifiers & 22.2 & locally enabled/disabled flags \\
\passthrough{\lstinline!Promise.try!} & 27.2.4.8 & normalize synchronous and async callback outcomes \\
\passthrough{\lstinline!Float16Array!}, related DataView/Math APIs & 21, 23, 25 & half-precision numeric storage and conversion \\
\end{longtable}

\begin{lstlisting}
// Iterator helpers: lazy map/filter then eager materialization
const evens = Iterator.from([1, 2, 3, 4, 5])
  .filter(n => n % 2 === 0).map(n => n * 10).toArray();
console.assert(evens.join(',') === '20,40');

const a = new Set([1, 2, 3]);
const b = new Set([3, 4]);
console.assert(a.intersection(b).has(3));
console.assert(a.union(b).size === 4);

const pattern = new RegExp(RegExp.escape('a+b'));
console.assert(pattern.test('literal a+b token'));
const p = Promise.try(() => 42);
p.then(v => console.assert(v === 42));
const half = new Float16Array([1.5, 2.25]);
console.assert(half[0] === 1.5);
\end{lstlisting}

\section{Reading the normative algorithms}\label{reading-the-normative-algorithms}

\textbf{{[}SPEC{]}} ECMA-262 uses specification records, abstract closures, internal slots written \passthrough{\lstinline![[Name]]!}, internal methods such as \passthrough{\lstinline![[Get]]!}, return completions, and convention markers such as \passthrough{\lstinline!?!} and \passthrough{\lstinline"!"} (clause 5). The notation expresses conditions of algorithm evaluation, not runtime JavaScript syntax.

\begin{enumerate}
\def\labelenumi{\arabic{enumi}.}
\tightlist
\item
  Locate the grammar production or built-in method subclause.
\item
  Identify preconditions, assertions, steps that may throw, and host hooks.
\item
  Follow references to abstract operations until the observable result is clear.
\item
  Distinguish invariant-preserving abstract machinery from externally reachable JavaScript behavior.
\item
  Record all alternate paths: normal result, throw, rejection, completion with \passthrough{\lstinline!return!} or \passthrough{\lstinline!break!}.
\end{enumerate}

\section{Portable engineering checklist}\label{portable-engineering-checklist}

\begin{itemize}
\tightlist
\item[$\square$]
  Fix the target edition to ECMA-262:2025 and record runtime engine versions.
\item[$\square$]
  Distinguish spec language behavior, host hooks, and web legacy features.
\item[$\square$]
  Validate types and conversions at public interfaces.
\item[$\square$]
  Handle empty collections, \passthrough{\lstinline!NaN!}, \passthrough{\lstinline!-0!}, BigInt, Unicode, nullish values.
\item[$\square$]
  Test iterator closure, abrupt completion, rejection, and module-loading failures.
\item[$\square$]
  Verify 2025 APIs in actual supported targets; provide fallback policy where necessary.
\item[$\square$]
  Review Annex B only when targeting browser legacy compatibility.
\end{itemize}

\section{Reference sources}\label{reference-sources}

\begin{itemize}
\tightlist
\item
  Ecma International, \emph{ECMA-262: ECMAScript 2025 Language Specification}, 16th edition, June 2025, supplied PDF; official edition: https://www.ecma-international.org/publications-and-standards/standards/ecma-262/ .
\item
  Live standard navigation (may describe later editions): https://tc39.es/ecma262/ . Always use the uploaded 2025 edition for clause mappings above.
\end{itemize}

\end{document}
